\documentclass[11pt,a4paper]{article}

\usepackage[UTF8]{ctex}
\usepackage[margin=2cm]{geometry}
\usepackage{xcolor}
\usepackage{listings}
\usepackage{booktabs}
\usepackage{longtable}
\usepackage{array}
\usepackage{tabularx}
\usepackage{hyperref}
\usepackage{titlesec}
\usepackage{tcolorbox}
\usepackage{fancyhdr}
\usepackage{parskip}
\usepackage{enumitem}

\definecolor{primaryblue}{HTML}{1F4E79}
\definecolor{softgray}{HTML}{F5F5F5}
\definecolor{codebg}{HTML}{F8F8F2}
\definecolor{codekey}{HTML}{0066CC}

\hypersetup{
  colorlinks=true,
  linkcolor=primaryblue,
  urlcolor=primaryblue,
}

\titleformat{\section}
  {\Large\bfseries\color{primaryblue}}
  {\thesection}{0.6em}{}
\titleformat{\subsection}
  {\large\bfseries\color{primaryblue}}
  {\thesubsection}{0.5em}{}
\titleformat{\subsubsection}
  {\normalsize\bfseries}
  {\thesubsubsection}{0.4em}{}

\lstset{
  basicstyle=\ttfamily\footnotesize,
  backgroundcolor=\color{codebg},
  breaklines=true,
  frame=single,
  framesep=4pt,
  framerule=0pt,
  rulecolor=\color{softgray},
  xleftmargin=5pt,
  xrightmargin=5pt,
  keywordstyle=\color{codekey}\bfseries,
  commentstyle=\color{gray}\itshape,
  stringstyle=\color{orange},
  showstringspaces=false,
  columns=flexible,
  language=Python,
}

\pagestyle{fancy}
\fancyhf{}
\fancyhead[L]{\footnotesize CausalVerify v11 Pipeline Walkthrough}
\fancyhead[R]{\footnotesize \thepage}

\title{\vspace{-1cm}
  \textbf{\Large CausalVerify v11 Execution Pipeline in Full Detail} \\[4pt]
  \textbf{\Large CausalVerify v11 
}
\author{Bilingual Technical Walkthrough \\ 
\date{}

\begin{document}

\maketitle
\thispagestyle{fancy}

\noindent The complete pipeline has 6 stages. Below I explain step by step \textbf{what we do, how we do it, and what we produce}.

\noindent 

\vspace{0.5em}
\hrule
\vspace{0.5em}

% ─────────────────────────────────────────────────────────
\section{Stage 0 — Data Preparation}

\textbf{Purpose}: establish an evaluable input (question + true answer). \\

\subsection{Exp A: Real Papers}

\textbf{Script}: \texttt{src/pipeline/expand\_benchmark.py} + \texttt{src/pipeline/generate\_paper\_jsons.py}

\textbf{What it does}:
\begin{enumerate}[leftmargin=1.5em]
\item Call the OpenAlex API to query papers from top economics journals. \\
      
\item Use GPT-4o to generate \texttt{research\_question} / \texttt{data\_description} / \texttt{institutional\_context} from the metadata. \\
\item Human review and extension $\rightarrow$ 262 paper JSON files. \\
      
\end{enumerate}

\textbf{Input}: 3-D allocation table of (economics domain $\times$ method $\times$ difficulty). \\


\textbf{Output}:

\begin{lstlisting}[language=]
experiments/exp_a/papers/paper_XX_did_easy.json    (262 files)
  Each contains:
    - paper_id, authors, year, journal
    - research_question     shown to LLM
    - data_description      shown to LLM
    - institutional_context shown to LLM
    - ground_truth:         hidden from LLM, used for scoring
        method_family    (DID/EVENT_STUDY/IV/RDD)
        conclusion_direction  (positive/negative)
        main_effect
        key_robustness
\end{lstlisting}

\subsection{Exp B: Synthetic DGP Scenarios}

\textbf{Script}: \texttt{src/pipeline/generate\_exp\_b\_scenarios.py} + \texttt{extend\_exp\_b\_to\_100.py}

\textbf{What it does}:
\begin{enumerate}[leftmargin=1.5em]
\item For 4 methods (DID/ES/IV/RDD), synthesize CSV data with NumPy so that the true $\beta^*$ is completely known. \\
      
\item Use GPT-4o to generate natural-language research questions (hiding the method name). \\
      GPT-4o 
\item 30 hand-written + 70 programmatically extended = 100 scenarios. \\
      30 
\end{enumerate}

\textbf{Key design}: \texttt{seed=42} fixed, fully reproducible.

\textbf{Output}:

\begin{lstlisting}[language=]
experiments/exp_b/scenarios/s01.json    (100 scenario files)
  - scenario_id, method_family, difficulty
  - research_question, data_description
  - data_file: "experiments/exp_b/data/sXX_data.csv"
  - dgp_truth:
      type:      "did" / "event" / "iv" / "rdd"
      effect:    -0.3      <- exact true effect
      direction: "negative"

experiments/exp_b/data/s01_data.csv     (100 CSV files)
  1600 rows (200 unit x 8 period) x 6 columns
\end{lstlisting}

\newpage

% ─────────────────────────────────────────────────────────
\section{Stage 1 — LLM Generation / LLM 

\textbf{Purpose}: have each LLM produce R code + research plan for each task. \\

\subsection{Exp A Calls / Exp A 

\textbf{Script}: \texttt{src/pipeline/run\_exp\_a.py}

\textbf{What it does} (for each paper $\times$ each model):

\begin{lstlisting}[language=Python]
# System prompt
"You are a financial economist... Propose an identification strategy..."

# User prompt (fill in three fields
"""
Research Question: {research_question}
Available Data: {data_description}
Institutional Context: {institutional_context}

Please provide:
  1. Identification Strategy
  2. Treatment and Control
  3. Key Identifying Assumptions
  4. Threats to Identification
  5. R Code Outline
  6. Expected Results
  7. Validity Checklist
"""

-> Call LLM API -> save
\end{lstlisting}

\textbf{Output}:

\begin{lstlisting}[language=]
experiments/exp_a/outputs/paper_XX_claude-opus-4-6.json  (1,572 files)
  Contains:
    - paper_id, model, run_at
    - llm_response:
        model, input_tokens, output_tokens
        content: complete markdown answer (with R code block)
        stop_reason
\end{lstlisting}

\textbf{Scale}: 262 papers $\times$ 6 models = \textbf{1,572 outputs} \\
\textbf{Time}: $\sim$6--12 hours (rate-limited)
\textbf{Cost}: $\sim\$80$--$150$

\subsection{Exp B Calls / Exp B 

\textbf{Script}: \texttt{src/pipeline/run\_exp\_b.py}

\textbf{Differences from Exp A}:
\begin{enumerate}[leftmargin=1.5em]
\item Provide the LLM with the \textbf{first 5 CSV rows + the file path}. \\
      
\item The prompt requires \textbf{complete, runnable R code} (not an outline). \\
\item Compress the prompt to 3 sections.
\end{enumerate}

\begin{lstlisting}[language=Python]
# User prompt template
"""
Research scenario: {title}
Research question: {research_question}
Data description: {data_description}
Data file path: experiments/exp_b/data/sXX_data.csv
Data preview (first 5 rows):
{pandas.to_string}
Column types: {dtypes}

Please implement the full analysis in 3 sections:
  1. Identification Strategy
  2. R Code (complete, self-contained)
  3. Expected Findings
"""
\end{lstlisting}

\textbf{Output}:

\begin{lstlisting}[language=]
experiments/exp_b/outputs/sXX_claude-opus-4-6.json   (600 files)
  Same structure as Exp A. Key difference: the R code in
  'content' references a real CSV path that can later be
  executed by L2b+.
\end{lstlisting}

\textbf{Scale}: 100 scenarios $\times$ 6 models = \textbf{600 outputs} \\
\textbf{Time}: $\sim$3--5 hours
\textbf{Cost}: $\sim\$40$--$80$

\newpage

% ─────────────────────────────────────────────────────────
\section{Stage 2 — External Verification (Pillar A) = L2b+}

\textbf{Purpose}: verify correctness via code execution + coefficient comparison. \textbf{This is the paper's Channel A.} \\
\textbf{

\textbf{Script}: \texttt{src/pipeline/score\_l2b\_plus.py}

\subsection{Step 2.1 — L1 Check / L1 

\begin{lstlisting}[language=Python]
l1 = bool(content.strip())    # Non-empty output?
\end{lstlisting}

\subsection{Step 2.2 — L2a: Extract R Code}

\begin{lstlisting}[language=Python]
# Regex match markdown code block
pattern = r"```(?:r|R)\n(.*?)\n```"
r_code = extract(content, pattern)
l2a = r_code is not None
\end{lstlisting}

\subsection{Step 2.3 — L2b: Actually Execute Code}

\begin{lstlisting}[language=Python]
# Write temp R script -> execute via Rscript -> 60s timeout
# 
with open("/tmp/tmpXXX.R", "w") as f:
    f.write(f"options(warn = -1)\nsuppressMessages({{\n{r_code}\n}})")

result = subprocess.run(
    ["Rscript", "--vanilla", "/tmp/tmpXXX.R"],
    capture_output=True, text=True, timeout=60
)
l2b = (result.returncode == 0)
stdout = result.stdout
\end{lstlisting}

\subsection{Step 2.4 — L2b+: Extract Coefficient + Compare}

\begin{lstlisting}[language=Python]
# Three-tier extraction strategy
# Tier 1: labeled pattern "DID coef: -0.2987"
# Tier 2: line-position pattern (regression table)
# Tier 3: fallback scan (lines containing "treat"/"did"/"post")

estimated = extract_coefficient(stdout, method=scenario["method_family"])

# Compare against DGP truth
true_effect = scenario["dgp_truth"]["effect"]   # -0.3
rel_error = abs(estimated - true_effect) / abs(true_effect)
l2b_plus = (rel_error < 0.5)    # +/- 50% tolerance
\end{lstlisting}

\textbf{Output}:

\begin{lstlisting}[language=]
experiments/exp_b/l2b_plus_scores.csv   (600 rows)
  Each row:
    scenario_id, method, model, model_slug
    true_effect, estimated, rel_error
    L1, L2a, L2b, L2b_plus
    err_excerpt (if failed)

experiments/exp_b/l2b_plus_summary.json
  { "Opus": { "L1": 100, "L2a": 99, "L2b": 94, "L2b_plus": 50 }, ... }
\end{lstlisting}

\textbf{Typical results}:

\begin{center}
\begin{tabular}{lccc}
\toprule
\textbf{Model} & \textbf{L2a} & \textbf{L2b} & \textbf{L2b+} \\
\midrule
Opus    & 99  & 94 & \textbf{50} \\
GPT-4o  & 100 & 78 & \textbf{25} \\
Sonnet  & 100 & 51 & \textbf{21} \\
o3      & 100 & 50 & \textbf{11} \\
Kimi    & 100 & 45 & \textbf{5}  \\
Gemini  & 57  & 32 & \textbf{13} \\
\bottomrule
\end{tabular}
\end{center}

\begin{tcolorbox}[colback=softgray, colframe=primaryblue, boxrule=0.5pt]
\textbf{What Channel A tells the paper} / \textbf{Channel A 

``External verification is feasible; L2b+ can precisely measure which model outputs are numerically correct.''

``
\end{tcolorbox}

\newpage

% ─────────────────────────────────────────────────────────
\section{Stage 3 — Internal Verification (Pillar B) = Calibration}

\textbf{Purpose}: measure whether the model's confidence in its own output matches actual correctness. \textbf{This is the paper's Channel B --- new in v11.} \\
\textbf{

\subsection{3.1 Pilot (just completed) / Pilot

\textbf{Script}: \texttt{src/pipeline/run\_calibration\_pilot.py}

\textbf{What it does}:
\begin{enumerate}[leftmargin=1.5em]
\item From \texttt{l2b\_plus\_scores.csv}, sample 10 L2b+ failures + 5 L2b+ passes. \\
      
\item For each case, re-query Opus with a retrospective prompt: \\
      
\end{enumerate}

\begin{lstlisting}[language=]
Here's your earlier analysis: {original_output}

Rate your confidence (0-1) on three dimensions:
  - method_confidence
  - specification_confidence
  - numerical_confidence

Output JSON only. Do NOT reveal true value.
\end{lstlisting}

\textbf{Output}:

\begin{lstlisting}[language=]
experiments/exp_b/calibration_pilot_results.csv   (15 rows)
  scenario_id, method, L2b_plus,
  method_confidence, specification_confidence, numerical_confidence
\end{lstlisting}

\textbf{Key numbers from the pilot / Pilot 

\begin{itemize}[leftmargin=1.5em]
\item Mean numerical\_conf on wrong cases = \textbf{0.50}
\item Mean numerical\_conf on correct cases = \textbf{0.61}
\item \textbf{Gap = +0.11} (supports ``confidently wrong'' hypothesis) \\
      \textbf{Gap = +0.11} (
\end{itemize}

\subsection{3.2 Full Experiment (next)}

\textbf{What it does}:
\begin{enumerate}[leftmargin=1.5em]
\item Run calibration elicitation for all 100 scenarios $\times$ 6 models. \\
      
\item Compute Expected Calibration Error (ECE).
\item Compute over-confidence ratio.
\end{enumerate}

\begin{lstlisting}[language=Python]
# Divide confidence into 10 bins
# For each bin:
#   actual pass rate = proportion of L2b+=1 in that bin
#   avg confidence = mean confidence within bin
# ECE = weighted average |actual pass rate - confidence|

# Over-confidence ratio:
P(conf > 0.8 | L2b+ = 0) / P(conf > 0.8 | L2b+ = 1)
\end{lstlisting}

\textbf{Expected output}:

\begin{lstlisting}[language=]
experiments/exp_b/calibration_scores.csv      (600 rows)
experiments/exp_b/calibration_summary.json    (6-model summary)
  Per model: ECE_method, ECE_specification, ECE_numerical
             P(conf>0.8 | L2b+=0), P(conf>0.8 | L2b+=1)
             over_confidence_ratio

paper/figures/calibration_reliability.pdf
  6 subplots, confidence bins x actual pass rate
  Diagonal = perfect calibration
  Curve deviating = miscalibration
\end{lstlisting}

\begin{tcolorbox}[colback=softgray, colframe=primaryblue, boxrule=0.5pt]
\textbf{What Channel B tells the paper} / \textbf{Channel B 

``Internal verification fails; self-reported confidence cannot be used as a triage signal.''

``
\end{tcolorbox}

\newpage

% ─────────────────────────────────────────────────────────
\section{Stage 4 — Cross-Channel Analysis}

\textbf{Purpose}: cross-compare results from Channel A and Channel B to derive the paper's core finding --- ``asymmetry.'' \\
\textbf{

\subsection{4.1 Head-to-head Ranking}

\textbf{Scripts}: \texttt{head\_to\_head\_ranking.py} + \texttt{bootstrap\_head\_to\_head.py}

\textbf{What it does}:
\begin{enumerate}[leftmargin=1.5em]
\item Use L2b+ pass rate as ``ground-truth ranking.''
\item Compute Kendall $\tau$ between each metric's ranking and the GT.
\item Bootstrap 10,000 times to get 95\% CI. / Bootstrap 10,000 
\end{enumerate}

\textbf{Output}:

\begin{lstlisting}[language=]
experiments/exp_b/head_to_head_ranking.json
  {
    "L2b vs GT":   { "tau": +0.73, "95CI": [+0.33, +0.87] },
    "L4_S1 vs GT": { "tau": -0.07, "95CI": [-0.41, +0.36] },
    "Calibration ECE vs GT": { ... }   <- NEW
  }
\end{lstlisting}

\subsection{4.2 Asymmetry Analysis (new in v11) / Asymmetry 

\textbf{What it does}: for each model, construct a ``channel profile'':

\begin{center}
\begin{tabular}{lcc}
\toprule
\textbf{Model} & \textbf{L2b+ pass rate} & \textbf{ECE} \\
\midrule
Opus    & 0.50 & 0.34 \\
GPT-4o  & 0.25 & 0.42 \\
...     & ...  & ...  \\
\bottomrule
\end{tabular}
\end{center}

Then analyze:
\begin{enumerate}[leftmargin=1.5em]
\item \textbf{Correlation} between Channel A and Channel B. / Channel A vs Channel B 
\item Whether any model is \textbf{good on Channel A but poor on Channel B} (= the paper's asymmetry claim). \\
      
\item Scatter / PCA visualization of the ``profile.''
\end{enumerate}

\textbf{Output}:

\begin{lstlisting}[language=]
paper/figures/channel_asymmetry.pdf
  Scatter plot (x = L2b+, y = 1 - ECE = calibration quality)
  Diagonal: "two channels agree"
  Off-diagonal points: asymmetry (your core finding)
  
\end{lstlisting}

\newpage

% ─────────────────────────────────────────────────────────
\section{Stage 5 — Figure Generation}

\textbf{Scripts}: \texttt{paper/figures/make\_fig*.py}

\textbf{Figures to produce (v11 checklist)}:

\begin{center}
\begin{tabular}{|c|l|l|}
\hline
\textbf{Figure} & \textbf{Script} & \textbf{Content} \\
\hline
Fig 1 & \texttt{make\_two\_dim\_framework.py}   & Two-dimensional framework (done) \\
Fig 2 & \texttt{make\_fig2\_headline.py}        & L2b+ cascade + model ranking \\
Fig 3 & \texttt{make\_fig3\_clean.py}           & Method-level dotplot \\
Fig 4 & \texttt{make\_fig4\_cascade.py}         & Execution funnel \\
Fig 5 & \texttt{make\_fig\_error\_cdf.py}       & Relative error CDF \\
Fig 6 & \texttt{make\_calibration\_reliability.py} & Reliability diagram (NEW) \\
Fig 7 & \texttt{make\_channel\_asymmetry.py}    & Two-channel scatter (NEW) \\
\hline
\end{tabular}
\end{center}

\textbf{Each script}:
\begin{itemize}[leftmargin=1.5em]
\item Reads from \texttt{experiments/exp\_b/*.csv} or \texttt{*.json}
\item Uses matplotlib + \texttt{palette.py}
\item Outputs \texttt{paper/figures/*.pdf} + \texttt{*.png}
\end{itemize}

% ─────────────────────────────────────────────────────────
\section{Stage 6 — Paper Compilation}

\subsection{6.1 LaTeX Source Editing / LaTeX 

\textbf{File}: \texttt{paper/latex/neurips\_v11.tex}

\textbf{Section changes in v11 / v11 

\begin{itemize}[leftmargin=1.5em]
\item \S2 renamed to \textbf{Two-Dimensional Evaluation Framework} (+1.5 pages new) \\
      \S2 
\item \S3 adds \textbf{\S3.2 Pillar B: Self-Assessment Calibration} (+0.5 pages new) \\
      \S3 
\item \S4 adds \textbf{\S4.3 Calibration Results + \S4.4 Cross-Channel Asymmetry} (+1.5 pages new) \\
      \S4 
\item \S6 adds \textbf{Practitioner Implications} (+1 page new) / \S6 
\end{itemize}

\subsection{6.2 Compilation}

\textbf{Docker approach / Docker 

\begin{lstlisting}[language=bash]
docker run --rm -v $(pwd)/paper/latex:/work -w /work \
    texlive/texlive:latest \
    bash -c "pdflatex neurips_v11.tex && bibtex neurips_v11 && \
             pdflatex neurips_v11.tex && pdflatex neurips_v11.tex"
\end{lstlisting}

\textbf{Output}: \texttt{paper/latex/neurips\_v11.pdf} ($\sim$15 pages main + appendix)

\newpage

% ─────────────────────────────────────────────────────────
\section{One-Page Pipeline Overview}

\begin{tcolorbox}[colback=softgray, colframe=primaryblue, boxrule=0.5pt]
\small
\begin{verbatim}
STAGE 0 - Data Preparation
  Exp A: OpenAlex -> 262 paper_XX.json
  Exp B: NumPy seed=42 -> 100 sXX.json + 100 CSVs + beta* known
                                |
                                v
STAGE 1 - LLM Generation / LLM 
  run_exp_a.py: 262 x 6 = 1,572 outputs     ( ~$100, 12h )
  run_exp_b.py: 100 x 6 = 600 outputs       ( ~$60,  5h )
                                |
                                v
STAGE 2 - Pillar A (External Verification) = L2b+
  score_l2b_plus.py: L1 -> L2a -> L2b -> L2b+
    -> l2b_plus_scores.csv + summary.json
    Key result: L2b+ pass rate 5-50%, tau(L2b+, GT) = +0.73
                                |
                                v
STAGE 3 - Pillar B (Internal Verification) = Calibration  [NEW]
  Pilot: 15 samples -> gap = 0.11 ('confidently wrong' holds)
  Full:  600 calibration queries
    -> calibration_scores.csv + summary.json
    Key result (expected): ECE > 0.3, overconfidence asymmetric
                                |
                                v
STAGE 4 - Cross-Channel Analysis
  head_to_head_ranking.py + bootstrap_head_to_head.py
  + asymmetry analysis (A vs B correlation per model)
    -> head_to_head_ranking.json + channel_asymmetry.pdf
                                |
                                v
STAGE 5 - Figure Generation
  Figures 1-7 (PDF + PNG)
                                |
                                v
STAGE 6 - Paper Compilation
  Docker + pdflatex -> neurips_v11.pdf ( ~15 pages )
\end{verbatim}
\end{tcolorbox}

\section{Each Stage's ``Paper Argument''}

\begin{center}
\begin{tabular}{|c|l|l|}
\hline
\textbf{Stage} & \textbf{Paper section} & \textbf{Argument carried} \\
\hline
0 & \S3.1 Methodology     & Dataset + ground truth construction \\
1 & \S4.1 Setup           & How 6 models are tested \\
2 & \S4.2 Pillar A        & Channel A is feasible ($\tau = +0.73$) \\
3 & \S4.3 Pillar B        & Channel B fails (ECE $>$ 0.3) \\
4 & \S4.4 Cross-Channel   & \textbf{Core finding: asymmetry} \\
5 & Figures               & Visual evidence \\
6 & \S6 Discussion        & Practitioner implications \\
\hline
\end{tabular}
\end{center}

\section{Where You Are Now}

\begin{tcolorbox}[colback=softgray, colframe=primaryblue, boxrule=0.5pt]
\begin{verbatim}
[DONE]    Stage 0 - Data preparation (v10 already done)
[DONE]    Stage 1 - LLM generation (v10 already done)
[DONE]    Stage 2 - Pillar A (v10 already done)
[PARTIAL] Stage 3 - Pillar B: pilot done, full pending   <- YOU ARE HERE
[TODO]    Stage 4 - Cross-Channel analysis
[TODO]    Stage 5 - New figures
[TODO]    Stage 6 - LaTeX editing + compilation
\end{verbatim}
\end{tcolorbox}

\section{Three Choices Ahead}

\begin{enumerate}[leftmargin=1.5em]
\item \textbf{Scale to the full calibration experiment} (Stage 3 complete, Day 2--4) --- \$30--60, a few hours. \\
      \textbf{
\item \textbf{First improve the prompt and script} (more rigorous contemporaneous elicitation) --- half a day. \\
      \textbf{
\item \textbf{First write \S2 Framework LaTeX} (Stage 6 ramp-up, no new data) --- 2 hours. \\
      \textbf{
\end{enumerate}

\end{document}